% HC-91 Emulator -- User Manual.  Build: make manual  (pdflatex x2)
\documentclass[11pt,a4paper]{article}
\usepackage[margin=2.6cm]{geometry}
\usepackage{parskip}
\usepackage{booktabs}
\usepackage{tabularx}
\usepackage{xcolor}
\usepackage{fancyhdr}
\usepackage{fancyvrb}
\usepackage[colorlinks=true,linkcolor=blue!50!black,urlcolor=blue!50!black]{hyperref}

\definecolor{shade}{gray}{0.95}
\fvset{frame=single,framerule=0.2pt,fontsize=\small,rulecolor=\color{black!30}}
\newcommand{\opt}[1]{\texttt{#1}}
\newcommand{\key}[1]{\textsf{\textbf{#1}}}

\pagestyle{fancy}
\fancyhf{}
\fancyhead[L]{\small HC-91 Emulator}
\fancyhead[R]{\small User Manual}
\fancyfoot[C]{\small\thepage}

\title{\Huge HC-91 Emulator\\[0.6ex]\LARGE User Manual}
\author{An emulator for the I.C.E.~Felix HC family\\
        (HC-85 / HC-90 / HC-91 / HC-128 / HC-2000)}
\date{June 2026}

\begin{document}
\maketitle
\thispagestyle{empty}

\begin{center}
\begin{tabular}{ll}
\toprule
CPU & Z80A (MMN80CPU) @ 3.5\,MHz, cycle-exact \\
Video & 256$\times$192, 15 colours, beam-accurate, floating bus \\
Sound & 1-bit beeper; AY-3-8912 on the HC-128 \\
Storage & cassette (.tap/.tzx, pulse-exact), floppy + CP/M on the HC-2000 \\
Build & plain C99, zero dependencies (\texttt{make}) \\
\bottomrule
\end{tabular}
\end{center}

\vspace{1em}
\tableofcontents
\clearpage

%=====================================================================
\section{Introduction}

\texttt{hc91emu} emulates the Romanian I.C.E.~Felix \emph{HC} home
computers --- ZX~Spectrum-compatible machines produced in Bucharest
from 1985 onward. The primary model is the \textbf{HC-91}, whose
genuine ROM differs from the Sinclair 48K ROM in only 50 bytes, making
it fully compatible with the Spectrum software catalogue. The
\textbf{HC-85} and \textbf{HC-90} are banner variants; the
\textbf{HC-128} adds 128K of banked RAM and an AY-3-8912 sound chip;
the \textbf{HC-2000} adds a floppy-disk interface and boots CP/M~2.2.

The emulator has two personalities:

\begin{itemize}
\item A \textbf{headless, scriptable} mode (the default): you state how
many 50\,Hz frames to run, schedule keypresses at given frames, and
collect results as PNG screenshots, OCR'd screen text, raw
framebuffers, WAV audio recordings, snapshots or input recordings.
Everything is deterministic, which makes it ideal for automated
testing --- the emulator's own test suite is built on it.
\item An \textbf{interactive SDL2} mode (\opt{--sdl}): a real window at
50\,Hz with live audio, host keyboard mapping, game-controller
support, turbo, pause, and cassette controls. The SDL library is
loaded at \emph{run time} with \texttt{dlopen}, so building the
emulator requires nothing beyond a C compiler.
\end{itemize}

Accuracy highlights: per-M-cycle memory/IO contention, beam-accurate
rendering (mid-frame border effects and per-scanline multicolour are
correct), the undocumented Z80 behaviours (MEMPTR, the Q register,
interrupted block instructions), exact tape pulse timing, and a
drift-free 69\,888-T-state frame. The core passes zexdoc, zexall,
Patrik Rak's \texttt{z80full}/\texttt{z80ccf}/\texttt{z80memptr}
(160/160 each) and all 162{,}000 SingleStepTests bus-level vectors.

%=====================================================================
\section{Quick start}

\begin{Verbatim}
make                 # build (no dependencies)
make test            # run the test suite

# boot to BASIC and look at the screen
./build/hc91emu --frames 250 --text

# play a game, interactively, with sound
./build/hc91emu --sdl game.tap --autoload
\end{Verbatim}

A single optional positional argument names the file to load,
dispatched by extension: \texttt{.tap}, \texttt{.tzx} (tape),
\texttt{.sna}, \texttt{.z80}, \texttt{.szx} (snapshots),
\texttt{.scr} (raw screen), \texttt{.rzx} (input recording).

%=====================================================================
\section{The emulated machines}

Select with \opt{--machine}; each machine has a sensible default ROM
from \texttt{roms/} (override with \opt{--rom}).

\begin{tabularx}{\textwidth}{lX}
\toprule
\textbf{Machine} & \textbf{Description} \\
\midrule
\opt{hc91} & Default. 48K Spectrum-compatible. Port \texttt{0x7E}
bit~0 pages a separate low 16K RAM bank over the ROM --- the CP/M mode
of the real machine, entered with \texttt{RANDOMIZE USR 14446}. \\
\opt{hc85}, \opt{hc90} & The earlier 48K models (banner-variant
genuine ROMs). \\
\opt{48k} & A standard Sinclair ZX Spectrum 48K. \\
\opt{hc128} & HC-91-derived 128K machine: RAM banks 0--7 at
\texttt{0xC000} via port \texttt{0x7FFD} (shadow screen in bank~7, ROM
select honoured with \opt{--rom1}, lock bit; odd banks contended) and
an AY-3-8912 at \texttt{0xFFFD}/\texttt{0xBFFD} mixed into the audio
output. \\
\opt{hc2000} & 48K-class machine plus the HC disk interface: i8272
floppy controller, IF1-style shadow ROM with 16K interface RAM, and a
system latch that can map the CP/M ROM and boot CP/M~2.2 from disk
(section~\ref{sec:disk}). \\
\bottomrule
\end{tabularx}

%=====================================================================
\section{Interactive use: the SDL frontend}

\begin{Verbatim}
./build/hc91emu --sdl game.tap --autoload
\end{Verbatim}

opens a resizable 640$\times$480 window running at 50\,Hz with live
beeper (and AY) audio; pacing is audio-clocked. The frontend needs
only the SDL2 \emph{runtime} library (\texttt{libSDL2-2.0.so.0}).

\subsection{Keyboard}

Letters, digits, \key{Enter} and \key{Space} map directly to the
Spectrum matrix.

\begin{tabularx}{\textwidth}{lX}
\toprule
\textbf{Host key} & \textbf{Spectrum meaning} \\
\midrule
\key{Shift} & CAPS SHIFT \\
\key{Ctrl} or \key{left Alt} & SYMBOL SHIFT \\
\key{Backspace} & DELETE (CAPS+0) \\
\key{Esc} & BREAK (CAPS+SPACE) \\
\key{arrows} & BASIC cursors (CAPS+5/6/7/8); plain 5/6/7/8 with
\opt{--joy-type cursor}; Kempston when the interface is attached \\
\key{, . ; ' - = /} & the corresponding SYMBOL SHIFT symbols \\
\bottomrule
\end{tabularx}

\subsection{Emulator controls}

\begin{tabularx}{\textwidth}{lX}
\toprule
\textbf{Key} & \textbf{Action} \\
\midrule
\key{Tab} & turbo while held (no pacing) \\
\key{F5} & pause / resume the emulation \\
\key{F6} & tape: manual play/stop \\
\key{F7} & tape: rewind (paused at the start) \\
\key{F8} & tape: swap cassette side (with \opt{--tape-b}; toggles) \\
\key{F10} / close & quit \\
\bottomrule
\end{tabularx}

\subsection{Game controllers}

The first SDL game controller drives the Kempston interface (d-pad,
left stick, and the A button as fire); plugging one in attaches the
interface automatically. On the keyboard, with \opt{--kempston},
arrows + \key{right Alt} drive Kempston instead of the cursor keys.

%=====================================================================
\section{Loading software from tape}

\subsection{The two loading paths}

\begin{description}
\item[Instant (default).] LOAD requests are satisfied by a ROM trap at
LD-BYTES: blocks appear immediately. Works for normal \texttt{.tap}
files and the standard/turbo data blocks of \texttt{.tzx} files.
\item[Pulse-level (\opt{--real-tape}).] The whole tape is compiled to
an EAR pulse stream with exact T-state timing and `played'. Required
for custom and turbo loaders, loading stripes, and anything
timing-sensitive. \opt{--play-at}~\textit{N} chooses when PLAY is
pressed (default: frame 320 with \opt{--autoload}).
\end{description}

\opt{--autoload} types \texttt{LOAD ""} + ENTER at frame 250 so tapes
start by themselves.

\subsection{Multi-load games and the cassette deck}
\label{sec:multiload}

Games that load levels on demand (``stop the tape'', ``turn the tape
over'') work in \opt{--real-tape} mode:

\begin{itemize}
\item The player \textbf{auto-pauses} at a block boundary whenever no
loader is actually polling the tape port. (Reads made by the ROM
keyboard scanner are recognised and ignored, so a ``press any key''
prompt does not keep the tape rolling past data.)
\item It \textbf{auto-resumes} as soon as loading restarts.
\item TZX \textbf{stop-the-tape} markers pause unconditionally.
\item At the \textbf{end of the tape}, if a loader is still searching,
the tape rewinds itself (``rewind tape'' prompts).
\item \textbf{Two-sided originals}: attach the other side with
\opt{--tape-b}~\textit{FILE} and press \key{F8} when the game asks you
to turn the tape over (headless: \opt{--swap-at}~\textit{N}). Each
swap toggles the sides, rewound and paused.
\end{itemize}

\begin{Verbatim}
./build/hc91emu --sdl shooter/p47_side_a.tzx --autoload --real-tape \
    --tape-b shooter/p47_side_b.tzx
# F8 at "PLEASE TURN TAPE OVER", any key to continue;
# hold Tab to fast-forward the long level searches
\end{Verbatim}

The tape state is reported on stderr (\texttt{tape: paused at block
5/23 ...}, \texttt{tape: resumed}, \texttt{tape: end of tape -
rewound}).

\subsection{Saving to tape}

\opt{--save-tape}~\textit{FILE} captures BASIC \texttt{SAVE} output
(via the SA-BYTES trap) as a byte-exact \texttt{.tap}.

\subsection{TZX block support}

Standard/turbo data (0x10/0x11), pure tone / pulse sequences / pure
data (0x12--0x14), \textbf{CSW recordings} (0x18; RLE and zlib Z-RLE),
\textbf{generalized data} (0x19; symbol alphabets, PRLE pilot runs,
polarity flags), pause/stop (0x20), groups (0x21/0x22), loops
(0x24/0x25), stop-if-48K (0x2A), and the information blocks
(0x30--0x35, 0x5A).

%=====================================================================
\section{The HC-128: banked RAM and AY sound}

The HC-128 keeps the 48K BASIC (its ROM is HC-91-derived) and adds the
Spectrum-128 hardware interfaces, so banking and sound are driven with
\texttt{OUT}s:

\begin{Verbatim}
# select RAM bank N at 0xC000 (keep bit 4 set): OUT 32765, 16+N
# move the stack down first or the switch pulls it from under BASIC!
CLEAR 32767: OUT 32765,17: POKE 49152,42: OUT 32765,16
\end{Verbatim}

The AY-3-8912 sits at the standard ports: \texttt{OUT 65533,reg}
selects, \texttt{OUT 49149,val} writes (reading \texttt{IN 65533}
returns the selected register). All three tone channels, noise and
envelopes are emulated with the measured volume curve and mixed
sample-accurately with the beeper into \opt{--wav} and the SDL audio.

128K \texttt{.z80} (hardware mode 3) and \texttt{.szx} (machine id 2)
snapshots save and load the full bank set, the \texttt{0x7FFD} latch
and the AY registers. 48K snapshots load into a locked, USR0-style
bank layout, so ordinary games run on the HC-128 too.

%=====================================================================
\section{The HC-2000: floppy disks and CP/M}
\label{sec:disk}

\subsection{Disk images}

\opt{--disk}~\textit{FILE} inserts drive A, \opt{--disk-b} drive B.
Raw images are recognised by size:

\begin{tabularx}{\textwidth}{llX}
\toprule
\textbf{Size} & \textbf{Geometry} & \textbf{Filesystem} \\
\midrule
655\,360 & 80$\times$2$\times$16$\times$256 & HC BASIC 3.5'' (640K) \\
737\,280 & 80$\times$2$\times$9$\times$512 & CP/M 2.2 3.5'' (720K, 4 boot tracks) \\
327\,680 & 40$\times$2$\times$16$\times$256 & HC BASIC 5.25'' \\
368\,640 & 40$\times$2$\times$9$\times$512 & CP/M 2.2 5.25'' \\
\bottomrule
\end{tabularx}

Writes are buffered in memory and written back to the image file when
the emulator exits.

\subsection{HC BASIC disks}

The interface presents Sinclair Interface~1 syntax to BASIC:

\begin{Verbatim}
CAT 1                          list the disk in drive 1
LOAD *"d";1;"run"              load a program from disk
\end{Verbatim}

(\texttt{CAT} is an extended-mode keyword: press CAPS+SYM, then
SYM+9.)

\subsection{Booting CP/M 2.2}

\begin{Verbatim}
./build/hc91emu --machine hc2000 --boot-cpm --disk cpm22-hc.img \
    --sdl
\end{Verbatim}

boots from the system tracks to \texttt{56k CP/M ver 2.2/2D} and the
\texttt{A>} prompt; \texttt{DIR}, applications (WordStar, Turbo
Pascal, MBASIC\dots) and the rest of CP/M work. The authentic route
from BASIC is \texttt{RANDOMIZE USR 14446}, which selects the CP/M
ROM and reboots into it. The CP/M console renders with its own font at
a relocated screen --- use screenshots rather than \opt{--text} there.

Set \texttt{HC91\_FDC\_DEBUG=1} to log every i8272 command/result on
stderr.

%=====================================================================
\section{Snapshots, screens and input recordings}

\subsection{Snapshots}

Load by passing the file; save with \opt{--save-sna},
\opt{--save-z80}, \opt{--save-szx} after the run.

\begin{tabularx}{\textwidth}{lX}
\toprule
\opt{.sna} & 48K, load + save. Save needs SP outside ROM (the PC is
pushed); on hc128 only the current 48K view can be represented. \\
\opt{.z80} & v1/v2/v3 load, incl.\ RLE compression and 128K files;
saves as v2 with uncompressed pages (48K or full 128K with
latch+AY). \\
\opt{.szx} & zx-state v1.4 load + save (48K and 128K). Compressed RAM
pages in files from other emulators are inflated by a built-in
DEFLATE decoder --- still zero external dependencies. \\
\opt{.scr} & raw 6912-byte screens; loading parks the CPU so the
picture stays. \\
\bottomrule
\end{tabularx}

A HALT state survives snapshots (PC is re-pointed at the HALT, or the
SZX HALTED flag is used).

\subsection{RZX input recordings}

\begin{Verbatim}
# record a session (embeds a .z80 of the starting state)
./build/hc91emu --frames 450 --type 'p2+3*4\n@260' --rzx-record calc.rzx
# replay: inputs come from the file, no key events needed
./build/hc91emu calc.rzx --frames 460 --text     # prints 14 again
\end{Verbatim}

Replay is bit-exact: each frame runs for the recorded number of R/M1
fetches and every port read returns the recorded byte. When the
recording ends, the machine continues running live. Compressed RZX
files from other emulators replay; ``protected'' (encrypted) ones are
rejected.

%=====================================================================
\section{Scripting and automation}

Typical recipe: schedule input, run N frames, collect output, assert.

\begin{Verbatim}
./build/hc91emu --frames 450 --type 'p2+3*4\n@260' --text \
    --screenshot out.png --wav out.wav --save-z80 out.z80
\end{Verbatim}

\subsection{Typing: \opt{--type} and \opt{--keys}}

\opt{--type "STRING@FRAME"} types a string with BASIC key mapping:
lowercase letters press the bare key (producing K-mode keywords at the
start of a statement), digits and common symbols work as expected,
\verb|\n| is ENTER. Keys are held for 6 frames with 6-frame gaps
($\approx$ 8 characters/second), so leave room between scheduled
items. \opt{--keys "FRAME:NAMES"} presses one raw chord, e.g.\
\texttt{"272:CAPS+SYM"}; names are A--Z, 0--9, ENTER, SPACE, CAPS,
SYM. Both repeat up to 32 times.

\textbf{Keyword cheat sheet} (48K BASIC keyword entry):

\begin{tabularx}{\textwidth}{lX}
\toprule
K-mode (statement start) & \texttt{p}=PRINT \quad \texttt{j}=LOAD
\quad \texttt{o}=POKE \quad \texttt{s}=SAVE \quad \texttt{t}=RANDOMIZE
\quad \texttt{x}=CLEAR \quad \texttt{f}=FOR \quad \texttt{i}=INPUT \\
EXTEND mode & press CAPS+SYM first (one \opt{--keys} entry), then:
plain \texttt{o}=PEEK, plain \texttt{i}=CODE, plain \texttt{l}=USR;
with SYM: \texttt{SYM+O}=OUT, \texttt{SYM+I}=IN, \texttt{SYM+Z}=BEEP,
\texttt{SYM+9}=CAT \\
\bottomrule
\end{tabularx}

Example --- OUT from BASIC (used by the HC-128 tests):

\begin{Verbatim}
--keys '260:CAPS+SYM' --keys '272:SYM+O' --type '65533,0\n@284'
\end{Verbatim}

\subsection{Output channels}

\begin{tabularx}{\textwidth}{lX}
\toprule
\opt{--text} & 24$\times$32 screen OCR against the ROM font (inverse
video recognised; unmatched cells print \texttt{?}) \\
\opt{--screenshot} & 320$\times$240 PNG of the final frame
(beam-painted) \\
\opt{--fb-dump} & the same frame as raw RGBA bytes \\
\opt{--wav} & 44.1\,kHz mono WAV of the whole run \\
\opt{--trace-frames} & frame/PC progress on stderr \\
\bottomrule
\end{tabularx}

Exit status is 0 on success, 1 on argument/file errors --- suitable
for shell test scripts (see \texttt{tests/run\_tests.sh} for several
hundred lines of worked examples).

%=====================================================================
\section{The debugger}

A scriptable monitor with a full Z80 disassembler (documented and
undocumented opcodes). Addresses/ports are hex (\texttt{\$} or
\texttt{0x} prefixes accepted); counts are decimal.

\subsection{Command-line flags}

\begin{tabularx}{\textwidth}{lX}
\toprule
\opt{--monitor} & stop in the monitor before the first instruction \\
\opt{--break ADDR} & PC breakpoint (repeatable) \\
\opt{--watch ADDR} & memory \emph{write} watchpoint \\
\opt{--rwatch ADDR} & memory \emph{read} watchpoint (fires on opcode
fetches too) \\
\opt{--pwatch PORT} & I/O watchpoint; values $\le$ FF match the low
byte, larger values match the full 16-bit port \\
\opt{--debug "c;c;q"} & scripted monitor commands, consumed one per
stop; falls back to stdin, then auto-continues \\
\opt{--trace FILE} & per-instruction trace (pre-execution state) \\
\bottomrule
\end{tabularx}

\subsection{Monitor commands}

\begin{tabularx}{\textwidth}{lX}
\toprule
\opt{r}, \opt{regs} & registers, decoded flags, frame-relative T-state,
frame number, total T \\
\opt{d}, \opt{dis} \textit{[ADDR [N]]} & disassemble (default PC, 8
instructions; \texttt{pc}/\texttt{sp}/\texttt{hl} accepted as ADDR) \\
\opt{m}, \opt{mem} \textit{ADDR [N]} & hex + ASCII dump (default 64
bytes) \\
\opt{s}, \opt{step} \textit{[N]} & execute N instructions, stop again \\
\opt{c}, \opt{cont} & continue \\
\opt{b} / \opt{w} / \opt{wr} / \opt{wp} \textit{[ADDR]} & toggle
breakpoint / write / read / port watch; bare = list \\
\opt{t}, \opt{trace} \textit{FILE\,$|$\,off} & tracing on/off \\
\opt{q}, \opt{quit} & exit immediately (skips end-of-run saves) \\
\opt{h}, \opt{help} & summary \\
\bottomrule
\end{tabularx}

An empty input line repeats the last command. Example session:

\begin{Verbatim}
$ ./build/hc91emu --frames 3 --break 11CB \
      --debug "regs;dis pc 3;mem 0000 16;step 3;regs;cont"
*** breakpoint at $11CB
11CB: 47           LD B,A
dbg> regs
PC=11CB AF=0044 BC=0000 DE=FFFF HL=0000 ... T=28/69888 frame=0
...
\end{Verbatim}

%=====================================================================
\section{The game library}

\texttt{tools/get\_library.sh} downloads 36 period classics from the
World of Spectrum archive into \texttt{software/library/}, classified
by genre (platform, arcade, isometric, shooter, puzzle, adventure,
sports), including 128K AY versions of Cybernoid and Tetris for the
HC-128 and the authentic two-sided P-47 cassettes.
\texttt{tools/verify\_library.sh} smoke-loads every title and stores
screenshots under \texttt{tests/out/library/}. The library README
carries the full index (title, year, publisher, file). The files are
copyrighted period software: the scripts make the library
reproducible, nothing under \texttt{software/} is committed.

%=====================================================================
\section{Accuracy, testing and the suite}

\texttt{make test} runs the suite: timing and contention unit tests,
the disassembler tests, zexdoc, boot/BASIC checks, the beam-renderer
and floating-bus acceptance tests, snapshot round-trips, the tape and
TZX matrix (including multi-load pause/resume), joysticks, CP/M
paging, the debugger, SingleStepTests vectors (162{,}000 bus-level
cases, fetched once by \texttt{tests/get\_vectors.sh}), the SDL
frontend under dummy drivers, the HC-128 and HC-2000 machines (CP/M
boots against a golden screen) and a library smoke test.
\texttt{make test-full} adds zexall; \texttt{RUN\_Z80TEST=1} re-runs
Rak's in-machine suites.

Environment variables: \texttt{HC91\_FDC\_DEBUG=1} (i8272 command
log), \texttt{HC91\_TAPE\_DEBUG=1} (tape block progression),
\texttt{SDL\_VIDEODRIVER=dummy SDL\_AUDIODRIVER=dummy} (headless SDL).

%=====================================================================
\section{Complete option reference}

\newcommand{\optrow}[2]{\opt{#1} & #2 \\}
\subsection*{General}
\begin{tabularx}{\textwidth}{p{4.6cm}X}
\toprule
\optrow{--machine M}{hc91 (default), hc85, hc90, 48k, hc128, hc2000}
\optrow{--rom FILE}{main ROM override (default per machine)}
\optrow{--rom1 FILE}{hc128 second ROM (default: copy of ROM 0)}
\optrow{--rom-if1 FILE}{hc2000 interface ROM (default roms/hc2ki1.rom)}
\optrow{--frames N}{frames to run headless (default 300; 50 = 1\,s)}
\optrow{--turbo}{accepted, no-op (headless never paces)}
\optrow{--no-floating-bus}{unattached ports read 0xFF}
\bottomrule
\end{tabularx}

\subsection*{Output}
\begin{tabularx}{\textwidth}{p{4.6cm}X}
\toprule
\optrow{--screenshot FILE}{PNG of the final frame (320$\times$240)}
\optrow{--text}{print the final screen as ROM-font OCR text}
\optrow{--fb-dump FILE}{final frame as raw RGBA}
\optrow{--wav FILE}{44.1\,kHz mono WAV of the run}
\optrow{--save-sna FILE}{save .sna after the run}
\optrow{--save-z80 FILE}{save .z80 v2 after the run}
\optrow{--save-szx FILE}{save .szx after the run}
\optrow{--save-scr FILE}{save the raw 6912-byte screen}
\bottomrule
\end{tabularx}

\subsection*{Input scripting}
\begin{tabularx}{\textwidth}{p{4.6cm}X}
\toprule
\optrow{--type "S@F"}{type string S from frame F (BASIC mapping, \texttt{\textbackslash n} = ENTER; repeatable)}
\optrow{--keys "F:NAMES"}{raw chord at frame F, e.g.\ 100:CAPS+1 (repeatable)}
\optrow{--autoload}{type \texttt{LOAD ""} at frame 250}
\optrow{--joy "F-G:DIRS"}{hold U/D/L/R/F over frames F..G (repeatable)}
\optrow{--joy-type T}{kempston (default) / sinclair / cursor}
\optrow{--kempston}{attach the Kempston interface (port 0x1F)}
\bottomrule
\end{tabularx}

\subsection*{Tape}
\begin{tabularx}{\textwidth}{p{4.6cm}X}
\toprule
\optrow{--real-tape}{pulse-level playback (multi-load auto pause/resume/rewind)}
\optrow{--play-at N}{press PLAY at frame N (default 320 with autoload)}
\optrow{--tape-b FILE}{the other cassette side (\opt{F8} / \opt{--swap-at} toggles)}
\optrow{--swap-at N}{insert \opt{--tape-b} at frame N (headless)}
\optrow{--save-tape FILE}{capture SAVE output as byte-exact .tap}
\bottomrule
\end{tabularx}

\subsection*{Disks and CP/M (hc2000)}
\begin{tabularx}{\textwidth}{p{4.6cm}X}
\toprule
\optrow{--disk FILE}{drive A image (raw .img/.dsk; written back on exit)}
\optrow{--disk-b FILE}{drive B image}
\optrow{--boot-cpm}{reset into the CP/M boot ROM}
\bottomrule
\end{tabularx}

\subsection*{Recordings}
\begin{tabularx}{\textwidth}{p{4.6cm}X}
\toprule
\optrow{--rzx-record FILE}{record the session as a replayable .rzx}
\bottomrule
\end{tabularx}

\subsection*{SDL frontend}
\begin{tabularx}{\textwidth}{p{4.6cm}X}
\toprule
\optrow{--sdl}{interactive window, 50\,Hz, live audio}
\optrow{--sdl-frames N}{auto-quit after N frames}
\optrow{--trace-frames}{frame/PC progress on stderr}
\bottomrule
\end{tabularx}

\subsection*{Debugger}
\begin{tabularx}{\textwidth}{p{4.6cm}X}
\toprule
\optrow{--monitor}{stop before the first instruction}
\optrow{--break ADDR}{PC breakpoint (hex; repeatable)}
\optrow{--watch ADDR}{memory write watchpoint}
\optrow{--rwatch ADDR}{memory read watchpoint (incl.\ fetches)}
\optrow{--pwatch PORT}{I/O watchpoint ($\le$FF = low byte)}
\optrow{--debug "C;C;.."}{scripted monitor commands}
\optrow{--trace FILE}{per-instruction trace to FILE}
\bottomrule
\end{tabularx}

%=====================================================================
\section{Troubleshooting}

\begin{description}
\item[A game freezes at its loading screen.] Try
\opt{--real-tape}: it probably uses a custom/turbo loader the instant
trap cannot serve.
\item[A multi-load game asked me to flip the tape and stopped.]
Use the two-sided release with \opt{--tape-b} and press \key{F8} at
the prompt; for single-file tapes the player pauses and resumes
automatically --- watch the \texttt{tape:} messages on stderr, and use
\key{F6}/\key{F7} for manual control. Hold \key{Tab} through long
level searches.
\item[\opt{--text} prints only \texttt{?} characters.] The program
draws with a custom font (e.g.\ the CP/M console); use
\opt{--screenshot} instead.
\item[A game ignores the joystick.] Check \opt{--joy-type}; for
Kempston games also pass \opt{--kempston} (the port is detached by
default so that floating-bus games work).
\item[CP/M shows a black screen in the first seconds.] Normal: the
boot ROM relocates the video memory and clears it before the banner
appears ($\sim$2500 frames to the \texttt{A>} prompt).
\item[Arkanoid freezes at round start.] You passed
\opt{--no-floating-bus}; the game needs the floating bus (default).
\end{description}

\vspace{2em}
\begin{center}
\small
See also: \texttt{man docs/hc91emu.1}, \texttt{README.md},
\texttt{ROADMAP.md} (development history),
\texttt{software/library/README.md} (game index).
\end{center}

\end{document}
